\documentclass[manuscript]{acmart}

\AtBeginDocument{%
  \providecommand\BibTeX{{%
    Bib\TeX}}}

\setcopyright{none}

\usepackage{array}
\usepackage{mathtools}
\usepackage{stmaryrd}
\usepackage{mathpartir}
\usepackage{listings}
\usepackage{xcolor}

\lstset{
  basicstyle=\ttfamily\small,
  breaklines=true,
  columns=fullflexible,
  keepspaces=true,
  frame=single,
  showstringspaces=false
}

\begin{document}

\title{Implementing Y\_T with Recursive Types}

\author{Runbang Yan}
\email{rbyan25@stu.pku.edu.cn}
\affiliation{%
  \institution{School of Computer Science, Peking University}
  \city{Beijing}
  \country{China}
}

\renewcommand{\shortauthors}{Yan}

\maketitle

\section{Direct Implementation of Y\_T}

The recursive type used on Slide 23 is:

\begin{lstlisting}
DivT = mu X. X -> T
\end{lstlisting}

Thus, after unfolding, a value of type \texttt{DivT} behaves like a function from \texttt{DivT} to \texttt{T}.

In MoonBit, this recursive type can be encoded using an enum:

\begin{lstlisting}
enum DivT[T] {
  MkDivT((DivT[T]) -> T)
}
\end{lstlisting}

The constructor \texttt{MkDivT} corresponds to the explicit \texttt{fold} operation.

\begin{lstlisting}
fn[T] fold_t(g : (DivT[T]) -> T) -> DivT[T] {
  MkDivT(g)
}
\end{lstlisting}

The \texttt{unfold} operation is implemented by pattern matching:

\begin{lstlisting}
fn[T] unfold_t(x : DivT[T]) -> (DivT[T]) -> T {
  match x {
    MkDivT(g) => g
  }
}
\end{lstlisting}

Therefore, the direct implementation of \texttt{Y\_T} is:

\begin{lstlisting}
fn[T] y_t(f : (T) -> T) -> T {
  let body = fn(x : DivT[T]) -> T {
    f(unfold_t(x)(x))
  }

  body(fold_t(body))
}
\end{lstlisting}

This corresponds to the term on Slide 23:

\begin{lstlisting}
Y_T =
  lambda f : T -> T.
    (lambda x : DivT.
       f ((unfold [DivT] x) x))
    (fold [DivT]
       (lambda x : DivT.
          f ((unfold [DivT] x) x)))
\end{lstlisting}

The type of the direct implementation is:

\begin{lstlisting}
y_t : ((T) -> T) -> T
\end{lstlisting}

The correspondence between the formal term and the MoonBit code is:

\begin{itemize}
  \item \texttt{DivT[T]} corresponds to \texttt{DivT}.
  \item \texttt{fold\_t} corresponds to \texttt{fold [DivT]}.
  \item \texttt{unfold\_t} corresponds to \texttt{unfold [DivT]}.
  \item \texttt{body} corresponds to the function:
\end{itemize}

\begin{lstlisting}
lambda x : DivT.
  f ((unfold [DivT] x) x)
\end{lstlisting}

\section{Whether the Direct Implementation Works}

The direct implementation is well-typed.

Inside \texttt{body}, the key expression is:

\begin{lstlisting}
f(unfold_t(x)(x))
\end{lstlisting}

If \texttt{x} has type \texttt{DivT}, then:

\begin{lstlisting}
unfold_t(x) : DivT -> T
\end{lstlisting}

Therefore:

\begin{lstlisting}
unfold_t(x)(x) : T
\end{lstlisting}

Since \texttt{f} has type:

\begin{lstlisting}
f : T -> T
\end{lstlisting}

we get:

\begin{lstlisting}
f(unfold_t(x)(x)) : T
\end{lstlisting}

So:

\begin{lstlisting}
body : DivT -> T
\end{lstlisting}

Then:

\begin{lstlisting}
fold_t(body) : DivT
\end{lstlisting}

Finally:

\begin{lstlisting}
body(fold_t(body)) : T
\end{lstlisting}

Thus the whole function has type:

\begin{lstlisting}
y_t : ((T) -> T) -> T
\end{lstlisting}

However, this direct implementation does not work as a practical fixed-point operator in MoonBit.

MoonBit is a strict, call-by-value language. Therefore, before calling a function, MoonBit evaluates its argument first.

The problematic expression is:

\begin{lstlisting}
f(unfold_t(x)(x))
\end{lstlisting}

Before calling \texttt{f}, MoonBit first evaluates:

\begin{lstlisting}
unfold_t(x)(x)
\end{lstlisting}

Suppose:

\begin{lstlisting}
x = fold_t(body)
\end{lstlisting}

Then:

\begin{lstlisting}
unfold_t(x) = body
\end{lstlisting}

So:

\begin{lstlisting}
unfold_t(x)(x) = body(x)
\end{lstlisting}

But by the definition of \texttt{body}:

\begin{lstlisting}
body(x) = f(unfold_t(x)(x))
\end{lstlisting}

Therefore, evaluating \texttt{unfold\_t(x)(x)} requires evaluating \texttt{unfold\_t(x)(x)} again. This causes infinite recursion.

For example, the following program diverges:

\begin{lstlisting}
fn main {
  let result = y_t(fn(_ : Int) -> Int {
    42
  })

  println(result)
}
\end{lstlisting}

Although the function ignores its argument and returns \texttt{42}, MoonBit still evaluates the argument before entering the function body. The argument is the infinitely unfolding self-application.

Therefore, the direct \texttt{Y\_T} is typable, but it does not work operationally as a useful fixed-point operator in MoonBit.

\section{Making It Work}

To make the fixed-point operator work in a strict language, the recursive self-application must be delayed.

Instead of defining a fixed point for an arbitrary type \texttt{T}, we define a fixed point for function types.

The useful type is:

\begin{lstlisting}
Y : ((A -> B) -> A -> B) -> A -> B
\end{lstlisting}

That is, given a function transformer:

\begin{lstlisting}
(A -> B) -> A -> B
\end{lstlisting}

we produce a recursive function:

\begin{lstlisting}
A -> B
\end{lstlisting}

The recursive type for the working version is:

\begin{lstlisting}
DivAB = mu X. X -> A -> B
\end{lstlisting}

In MoonBit, this can be encoded as:

\begin{lstlisting}
enum DivAB[A, B] {
  MkDivAB((DivAB[A, B]) -> (A) -> B)
}
\end{lstlisting}

The explicit \texttt{fold} operation is:

\begin{lstlisting}
fn[A, B] fold_ab(g : (DivAB[A, B]) -> (A) -> B) -> DivAB[A, B] {
  MkDivAB(g)
}
\end{lstlisting}

The explicit \texttt{unfold} operation is:

\begin{lstlisting}
fn[A, B] unfold_ab(x : DivAB[A, B]) -> (DivAB[A, B]) -> (A) -> B {
  match x {
    MkDivAB(g) => g
  }
}
\end{lstlisting}

The working fixed-point operator is:

\begin{lstlisting}
fn[A, B] y(f : ((A) -> B) -> (A) -> B) -> (A) -> B {
  let body = fn(x : DivAB[A, B]) -> (A) -> B {
    f(fn(a : A) -> B {
      unfold_ab(x)(x)(a)
    })
  }

  body(fold_ab(body))
}
\end{lstlisting}

The important part is the additional function abstraction:

\begin{lstlisting}
fn(a : A) -> B {
  unfold_ab(x)(x)(a)
}
\end{lstlisting}

This lambda delays the recursive call.

In the direct version, the self-application is evaluated immediately:

\begin{lstlisting}
unfold_t(x)(x)
\end{lstlisting}

In the working version, the self-application is placed under a function:

\begin{lstlisting}
fn(a : A) -> B {
  unfold_ab(x)(x)(a)
}
\end{lstlisting}

Therefore, recursive unfolding only happens when the resulting function is applied to an actual argument.

\section{Complete MoonBit Implementation}

\begin{lstlisting}
enum DivAB[A, B] {
  MkDivAB((DivAB[A, B]) -> (A) -> B)
}

fn[A, B] fold_ab(g : (DivAB[A, B]) -> (A) -> B) -> DivAB[A, B] {
  MkDivAB(g)
}

fn[A, B] unfold_ab(x : DivAB[A, B]) -> (DivAB[A, B]) -> (A) -> B {
  match x {
    MkDivAB(g) => g
  }
}

fn[A, B] y(f : ((A) -> B) -> (A) -> B) -> (A) -> B {
  let body = fn(x : DivAB[A, B]) -> (A) -> B {
    f(fn(a : A) -> B {
      unfold_ab(x)(x)(a)
    })
  }

  body(fold_ab(body))
}
\end{lstlisting}

\section{Three Recursive Functions Defined Using the Fixed-Point Operator}

\subsection{Factorial}

\begin{lstlisting}
let fact : (Int) -> Int =
  y(fn(self : (Int) -> Int) -> (Int) -> Int {
    fn(n : Int) -> Int {
      if n <= 1 {
        1
      } else {
        n * self(n - 1)
      }
    }
  })
\end{lstlisting}

Example call:

\begin{lstlisting}
println("fact(5) = \{fact(5)}")
\end{lstlisting}

Expected result:

\begin{lstlisting}
fact(5) = 120
\end{lstlisting}

\subsection{Fibonacci}

\begin{lstlisting}
let fib : (Int) -> Int =
  y(fn(self : (Int) -> Int) -> (Int) -> Int {
    fn(n : Int) -> Int {
      if n <= 1 {
        n
      } else {
        self(n - 1) + self(n - 2)
      }
    }
  })
\end{lstlisting}

Example call:

\begin{lstlisting}
println("fib(10) = \{fib(10)}")
\end{lstlisting}

Expected result:

\begin{lstlisting}
fib(10) = 55
\end{lstlisting}

\subsection{Summation}

\begin{lstlisting}
let sum_to : (Int) -> Int =
  y(fn(self : (Int) -> Int) -> (Int) -> Int {
    fn(n : Int) -> Int {
      if n <= 0 {
        0
      } else {
        n + self(n - 1)
      }
    }
  })
\end{lstlisting}

Example call:

\begin{lstlisting}
println("sum_to(100) = \{sum_to(100)}")
\end{lstlisting}

Expected result:

\begin{lstlisting}
sum_to(100) = 5050
\end{lstlisting}

\section{Full MoonBit Program}

\begin{lstlisting}
enum DivAB[A, B] {
  MkDivAB((DivAB[A, B]) -> (A) -> B)
}

fn[A, B] fold_ab(g : (DivAB[A, B]) -> (A) -> B) -> DivAB[A, B] {
  MkDivAB(g)
}

fn[A, B] unfold_ab(x : DivAB[A, B]) -> (DivAB[A, B]) -> (A) -> B {
  match x {
    MkDivAB(g) => g
  }
}

fn[A, B] y(f : ((A) -> B) -> (A) -> B) -> (A) -> B {
  let body = fn(x : DivAB[A, B]) -> (A) -> B {
    f(fn(a : A) -> B {
      unfold_ab(x)(x)(a)
    })
  }

  body(fold_ab(body))
}

fn main {
  let fact : (Int) -> Int =
    y(fn(self : (Int) -> Int) -> (Int) -> Int {
      fn(n : Int) -> Int {
        if n <= 1 {
          1
        } else {
          n * self(n - 1)
        }
      }
    })

  let fib : (Int) -> Int =
    y(fn(self : (Int) -> Int) -> (Int) -> Int {
      fn(n : Int) -> Int {
        if n <= 1 {
          n
        } else {
          self(n - 1) + self(n - 2)
        }
      }
    })

  let sum_to : (Int) -> Int =
    y(fn(self : (Int) -> Int) -> (Int) -> Int {
      fn(n : Int) -> Int {
        if n <= 0 {
          0
        } else {
          n + self(n - 1)
        }
      }
    })

  println("fact(5) = \{fact(5)}")
  println("fib(10) = \{fib(10)}")
  println("sum_to(100) = \{sum_to(100)}")
}
\end{lstlisting}

Expected output:

\begin{lstlisting}
fact(5) = 120
fib(10) = 55
sum_to(100) = 5050
\end{lstlisting}

\section{Explicit Formulation with Fold and Unfold}

\subsection{Direct Version}

The recursive type is:

\begin{lstlisting}
DivT = mu X. X -> T
\end{lstlisting}

The explicit fold and unfold operations have the following types:

\begin{lstlisting}
fold [DivT]   : (DivT -> T) -> DivT
unfold [DivT] : DivT -> DivT -> T
\end{lstlisting}

The direct version is:

\begin{lstlisting}
Y_T =
  lambda f : T -> T.
    (lambda x : DivT.
       f ((unfold [DivT] x) x))
    (fold [DivT]
       (lambda x : DivT.
          f ((unfold [DivT] x) x)))
\end{lstlisting}

Its type is:

\begin{lstlisting}
Y_T : (T -> T) -> T
\end{lstlisting}

This version does not work under call-by-value evaluation because the expression:

\begin{lstlisting}
(unfold [DivT] x) x
\end{lstlisting}

must be evaluated before \texttt{f} is called. This causes infinite unfolding.

\subsection{Working Delayed Version}

The recursive type for the working version is:

\begin{lstlisting}
DivAB = mu X. X -> A -> B
\end{lstlisting}

The explicit fold and unfold operations have the following types:

\begin{lstlisting}
fold [DivAB]   : (DivAB -> A -> B) -> DivAB
unfold [DivAB] : DivAB -> DivAB -> A -> B
\end{lstlisting}

The working version is:

\begin{lstlisting}
Y =
  lambda f : (A -> B) -> A -> B.
    (lambda x : DivAB.
       f (lambda a : A.
            ((unfold [DivAB] x) x) a))
    (fold [DivAB]
       (lambda x : DivAB.
          f (lambda a : A.
               ((unfold [DivAB] x) x) a)))
\end{lstlisting}

Its type is:

\begin{lstlisting}
Y : ((A -> B) -> A -> B) -> A -> B
\end{lstlisting}

The expression:

\begin{lstlisting}
lambda a : A.
  ((unfold [DivAB] x) x) a
\end{lstlisting}

is the delay. It prevents the self-application from being evaluated immediately.

\section{Why the Delayed Version Is a Fixed-Point Operator}

Let:

\begin{lstlisting}
body =
  lambda x : DivAB.
    f (lambda a : A.
         ((unfold [DivAB] x) x) a)
\end{lstlisting}

Then:

\begin{lstlisting}
Y f
=
body (fold [DivAB] body)
\end{lstlisting}

Expanding \texttt{body}:

\begin{lstlisting}
Y f
=
f (lambda a : A.
     ((unfold [DivAB] (fold [DivAB] body))
       (fold [DivAB] body)) a)
\end{lstlisting}

By the iso-recursive type rule:

\begin{lstlisting}
unfold [DivAB] (fold [DivAB] body)
=
body
\end{lstlisting}

Therefore:

\begin{lstlisting}
Y f
=
f (lambda a : A.
     (body (fold [DivAB] body)) a)
\end{lstlisting}

Since:

\begin{lstlisting}
body (fold [DivAB] body)
=
Y f
\end{lstlisting}

we get:

\begin{lstlisting}
Y f
=
f (lambda a : A. (Y f) a)
\end{lstlisting}

By eta-equivalence for functions:

\begin{lstlisting}
lambda a : A. (Y f) a
=
Y f
\end{lstlisting}

Therefore:

\begin{lstlisting}
Y f = f (Y f)
\end{lstlisting}

So the delayed version behaves as a fixed-point operator for recursive functions.

\end{document}
\endinput
%%
%% End of file.
